\documentclass[11pt]{article}
\usepackage[T1]{fontenc}
\usepackage[utf8]{inputenc}
\usepackage{lmodern}
\usepackage[margin=0.92in]{geometry}
\usepackage{microtype}
\usepackage{setspace}
\usepackage{amsmath,amssymb,amsthm,mathtools,mathrsfs}
\usepackage{booktabs,array,tabularx,longtable}
\usepackage{enumitem}
\usepackage{xcolor}
\usepackage{hyperref}
\usepackage[nameinlink,noabbrev]{cleveref}
\usepackage{listings}
\usepackage{tikz}
\usetikzlibrary{arrows.meta,positioning,calc}
\hypersetup{colorlinks=true,linkcolor=blue!45!black,citecolor=green!35!black,urlcolor=blue!55!black,pdfauthor={OpenAI GPT-5.6 Pro},pdftitle={Catalan Cyclic Resultants and the Alaoglu--Erdos Conjecture}}
\setstretch{1.045}
\setlist[itemize]{leftmargin=1.5em,itemsep=2pt,topsep=4pt}
\setlist[enumerate]{leftmargin=1.7em,itemsep=2pt,topsep=4pt}
\allowdisplaybreaks

\newtheorem{theorem}{Theorem}[section]
\newtheorem{proposition}[theorem]{Proposition}
\newtheorem{lemma}[theorem]{Lemma}
\newtheorem{corollary}[theorem]{Corollary}
\newtheorem{criterion}[theorem]{Finishing criterion}
\theoremstyle{definition}
\newtheorem{definition}[theorem]{Definition}
\newtheorem{problem}[theorem]{Target problem}
\newtheorem{computation}[theorem]{Computational check}
\theoremstyle{remark}
\newtheorem{remark}[theorem]{Remark}

\DeclareMathOperator{\Res}{Res}
\DeclareMathOperator{\Tr}{Tr}
\DeclareMathOperator{\ord}{ord}
\DeclareMathOperator{\Gal}{Gal}
\DeclareMathOperator{\gcdpoly}{gcd}
\newcommand{\ZZ}{\mathbb Z}
\newcommand{\QQ}{\mathbb Q}
\newcommand{\RR}{\mathbb R}
\newcommand{\CC}{\mathbb C}
\newcommand{\FF}{\mathbb F}
\newcommand{\eps}{\varepsilon}
\newcommand{\zR}{\zeta_{\!\mathrm R}}
\newcommand{\statusproved}{\textnormal{\sffamily[proved here]}}
\newcommand{\statuscomputed}{\textnormal{\sffamily[checked symbolically/numerically]}}
\newcommand{\statusconditional}{\textnormal{\sffamily[conditional target]}}
\newcommand{\statusinherited}{\textnormal{\sffamily[inherited]}}


\newcommand{\GammaOne}{0.73693748}
\newcommand{\GammaTwo}{0.42547105}
\newcommand{\GammaThree}{0.30085347}
\newcommand{\GammaEight}{0.12282291}
\newcommand{\KappaThirteen}{1.17302149}
\newcommand{\KappaTwoEight}{0.72811919}
\newcommand{\KappaBalanced}{0.44490230}
\newcommand{\BalancedConstant}{0.73562416}


\title{\bfseries Catalan Cyclic Resultants, Diagonal Kummer Descent,\\
and an Adelic Tangent Lattice\\[4pt]
\large Twenty-Eighth Progress Report on the Alaoglu--Erd\H{o}s Conjecture}
\author{OpenAI GPT-5.6 Pro\\with the supplied unified report as inherited baseline}
\date{22 August 2026}

\begin{document}
\maketitle

\begin{abstract}
Assume, for contradiction, that a least nonintegral exponent \(\beta\) exists with
\(M=2^\beta\in\ZZ\) and \(A=3^\beta\in\ZZ\), as in the supplied unified report.
This continuation does not repeat that report's twenty-seven component reports.  It develops a new construction from the four adjacent \(\{2,3\}\)-smooth pairs
\[
(1,2),\qquad(2,3),\qquad(3,4),\qquad(8,9).
\]
For a lower convergent \(N/q<\beta\), each pair gives a small distinguished algebraic integer
\(((a+1)^{\beta q/N}-a^{\beta q/N})^N-1\).  Raising the difference to the same exponent \(N\) produces a \emph{diagonal Kummer descent}: although the two radicals generically generate a rank-two Kummer extension, the powered difference lies in a cyclic extension of degree at most \(N\).  Its conjugates are governed by a universal monic polynomial \(\Phi_{N;P,Q}\in\ZZ[T]\), expressible as a sparse resultant.

The principal new analytic result is a sharp boundary-layer law.  The full cyclic resultant has a main term of order \(\exp(N^2\log(a+1))\), but its non-distinguished correction is
\(\exp(-\gamma_a\sqrt N+o(\sqrt N))\), where
\[
\gamma_a=\frac{\zR(3/2)}{\sqrt{2\pi a(a+1)}}.
\]
The four resultants admit a unique primitive multiplicative combination in which both the entire \(N^2\)-height and the continued-fraction error cancel.  The resulting positive rational number satisfies
\[
\Theta_{N,q}
 =\frac{\mathscr R_1(N,q)^2\mathscr R_8(N,q)}
 {\mathscr R_3(N,q)\mathscr R_2(N,q)^2}
 =C_0\exp(-\kappa\sqrt N+o(\sqrt N)),
\qquad \kappa=\KappaBalanced\ldots>0.
\]
Consequently, any counterexample forces a stretched-exponential primitive denominator in this quotient.  This gives a concrete, finite, resultant--gcd criterion whose improvement by any fixed amount proves the conjecture.

At a prime numerator \(N=\ell\), the universal polynomial has an exact Frobenius collapse
\(\Phi_{\ell;P,Q}(T)\equiv(T-Q+P)^\ell\pmod\ell\).  Under the natural resonance \(Q-P\equiv1\pmod\ell\), the resultant is divisible by \(\ell^\ell\), and the first quotient is controlled by the truncated logarithm
\(H_\ell(X)=((X-1)^\ell-X^\ell+1)/\ell\).  When the two basic outputs satisfy
\(M^q\equiv2\) and \(A^q\equiv3\pmod\ell\), the four first Frobenius quotients obey exactly the same two integer tangent relations as the four real Catalan defects.  Their balanced residual quotient reduces to a one-variable rational map
\(6(t-1)/(t^2(3-t))\) over \(\FF_\ell\).

These results do not yet prove the Alaoglu--Erd\H{o}s conjecture.  They replace a diffuse ``control the other conjugates'' obstruction by two sharply stated arithmetic tasks: control a specific cross-resultant gcd at the \(e^{\kappa\sqrt N}\) scale, or derive enough repeated finite-place resonance to force that control.  All exact algebraic and congruence statements are suitable for direct Lean formalization.
\end{abstract}

\tableofcontents

\section{Status, scope, and external reconnaissance}

\subsection{Mathematical status}
The conclusion of this report is deliberately separated into three levels.
\begin{itemize}
  \item Every identity, resultant formula, tangent-lattice calculation, Frobenius congruence, and finite-field reduction marked \statusproved\ is proved in the text.
  \item The \(\sqrt N\)-boundary law is proved by a local expansion and a dominated Riemann-sum argument; the numerical table is only an independent check, not evidence substituted for proof.
  \item The final cross-resultant denominator estimate is a necessary condition for a counterexample and a sufficient target for a proof.  No unproved divisibility is silently assumed.
\end{itemize}
Thus the full conjecture remains open in this report.  The advance is a new exact interface at which a future proof can land.

\subsection{What was checked about the reported rival results}
Public searches made on 22 August 2026 do confirm that several of the reported developments have concrete public artifacts.  The ICARM Elliptic Curve Rank Leaderboard records curve~\#273 with thirty independent rational points and credits Claude, Levent Alp\H{o}ge, and Ava Howell; this verifies rank at least~\(30\), while the stated upper bound is conditional on GRH and BSD \cite{ICARM30}.  A public Hadamard construction repository reports hundreds of machine-verified orders and says that, below~\(2000\), only the construction files for orders~\(1212\) and~\(1940\) remain inaccessible even though existence is known \cite{HadamardRepo}; a separate public ``by Claude note from the Fable~5 era on order~\(668\) explicitly proves only a restricted-symmetry obstruction, not a full nonexistence or construction theorem \cite{Fable668}.  A 2026 arXiv paper also cites the announced Jacobian counterexample as an input \cite{JacobianMoment}.

I did not locate a publicly inspectable Alaoglu--Erd\H{o}s manuscript or Lean repository as of that date.  No external announcement is used as a lemma below: the argument is self-contained modulo the inherited facts explicitly identified in \cref{sec:baseline}.  The distinction matters because a leaderboard witness, a repository artifact, a preprint, a news report, and a kernel-checked proof have different trust boundaries.

\section{Inherited baseline and the non-duplication boundary}\label{sec:baseline}

The supplied unified report already establishes and organizes the following conditional framework.  We use it without reproducing its proofs.

\begin{enumerate}
\item If a nonintegral solution exists, there is a least positive nonintegral solution \(\beta\), and with
\[
M=2^\beta,\qquad A=3^\beta,
\]
all nonnegative solutions form the additive monoid
\(\ZZ_{\ge0}+\ZZ_{\ge0}\beta\).  In particular every \(\{2,3\}\)-smooth positive integer \(r\) has \(r^\beta\in\ZZ\). \statusinherited

\item The primitive pair \((M,A)\) satisfies strong perfect-power and valuation restrictions, \(\beta\) is transcendental, and extensive finite verification excludes a substantial initial interval. \statusinherited

\item The report has already examined general Kummer norms, radical-field product formulas, Pad\'e constructions, moment determinants, phase averaging, \(p\)-adic linearization, canonical products, orbit codings, and many denominator-aware variants.  In particular, it explicitly identifies the generic ``small distinguished conjugate versus enormous house'' failure. \statusinherited
\end{enumerate}

The present report adds neither another generic norm nor another analogy.  It exploits the exceptional fact that the \(\{2,3\}\)-smooth numbers contain exactly four adjacent pairs relevant here, and that each adjacent difference equals the unit \(1\).  This removes the exponential scale from the distinguished embedding, while a diagonal Kummer symmetry removes one radical direction from the powered defect.  The resulting four exact cyclic resultants can then be balanced against one another.

Throughout, assume a hypothetical least counterexample \(\beta\) and write
\[
\mathcal A=\{1,2,3,8\}.
\]
For \(a\in\mathcal A\), both \(a\) and \(a+1\) are \(\{2,3\}\)-smooth, hence
\[
B_a=a^\beta\in\ZZ,\qquad C_a=(a+1)^\beta\in\ZZ.
\]
The notation \(B_a,C_a\) is local to this report and should not be confused with asymptotic constants introduced later.

\section{The adjacent-smooth ladder and rational approach to the unit gap}

Let \(N/q<\beta\) be a reduced lower convergent of the continued fraction of \(\beta\), and put
\[
\eps=q\beta-N>0,\qquad s=\frac{q\beta}{N}=1+\frac{\eps}{N}.
\]
For lower convergents,
\[
0<\eps<\frac1q=O_\beta(N^{-1}),\qquad s-1=O_\beta(N^{-2}).
\]
For each \(a\in\mathcal A\), define
\[
P_a=B_a^q=a^{q\beta},\qquad Q_a=C_a^q=(a+1)^{q\beta},
\]
\[
u_a=P_a^{1/N}=a^s,\qquad v_a=Q_a^{1/N}=(a+1)^s,
\]
where the positive real roots are used at the distinguished embedding.  The real gap
\[
h_a(s)=v_a-u_a=(a+1)^s-a^s
\]
satisfies \(h_a(1)=1\).  Since
\[
h_a'(s)=(a+1)^s\log(a+1)-a^s\log a>0\quad(s\ge1),
\]
we have \(h_a(s)>1\) for every lower approximant.

\begin{definition}[distinguished Catalan defect]
For \(a\in\mathcal A\), set
\[
\delta_a(N,q)=h_a(s)^N-1
 =\left((a+1)^{q\beta/N}-a^{q\beta/N}\right)^N-1.
\]
\end{definition}

Although \(u_a\) and \(v_a\) may jointly generate a degree of order \(N^2\), the powered difference \((v_a-u_a)^N\) descends to one Kummer direction.  That is the first key point.

\section{Diagonal Kummer descent and the universal polynomial}\label{sec:kummer}

\subsection{The degree drop}
Fix positive integers \(P<Q\) and an integer \(N\ge2\).  Let
\[
u=P^{1/N},\qquad v=Q^{1/N},\qquad \alpha=\frac vu,
\]
and let \(\zeta_N=e^{2\pi i/N}\).  Put
\[
Z_j=(\zeta_N^jv-u)^N=P(\zeta_N^j\alpha-1)^N,
\qquad 0\le j<N.
\]
The separate radicals \(u,v\) can require two Kummer generators.  However, every \(Z_j\) belongs to
\[
\QQ(\alpha),\qquad \alpha^N=\frac QP,
\]
whose degree is at most \(N\).  The simultaneous scalar action
\((u,v)\mapsto(\zeta_Nu,\zeta_Nv)\) disappears after taking the \(N\)-th power.  We call this \emph{diagonal Kummer descent}.

\begin{theorem}[universal diagonal polynomial]\label{thm:universal-poly}
Define
\[
\Phi_{N;P,Q}(T)=\prod_{j=0}^{N-1}\left(T-P(\zeta_N^j\alpha-1)^N\right).
\]
Then \(\Phi_{N;P,Q}(T)\in\ZZ[T]\) is monic of degree \(N\), possibly with repeated roots, and
\begin{equation}\label{eq:resultant}
\Phi_{N;P,Q}(T)
=P^{-N}\Res_X\!\left(PX^N-Q,\;T-P(X-1)^N\right).
\end{equation}
Its roots are precisely the \(Z_j\).
\end{theorem}

\begin{proof}
Each \(\zeta_N^jv-u\) is an algebraic integer, hence so is \(Z_j\).  The multiset \(\{Z_j\}\) is stable under every automorphism of a splitting field of \(X^N-Q/P\); therefore every coefficient of the displayed product is both rational and an algebraic integer, hence an ordinary integer.  For the resultant identity, the roots of \(PX^N-Q\) are \(\zeta_N^j\alpha\), and the standard root formula for a resultant gives
\[
\Res_X(PX^N-Q,T-P(X-1)^N)
=P^N\prod_j\left(T-P(\zeta_N^j\alpha-1)^N\right).
\]
\end{proof}

\begin{definition}[cyclic Catalan resultant]
Set
\begin{equation}\label{eq:cyclicR}
\mathscr R_N(P,Q)
=\prod_{j=0}^{N-1}(Z_j-1)
=(-1)^N\Phi_{N;P,Q}(1)\in\ZZ.
\end{equation}
For the pair \((a,a+1)\) and the lower convergent \(N/q\), write
\[
\mathscr R_a(N,q)=\mathscr R_N(P_a,Q_a).
\]
\end{definition}

\begin{proposition}[positivity and nonvanishing]\label{prop:positive}
For the lower-convergent data above, \(\mathscr R_a(N,q)>0\).
\end{proposition}

\begin{proof}
The distinguished factor is \(h_a(s)^N-1>0\).  Factors indexed by \(j\) and \(N-j\) are complex conjugates and have positive product.  If \(N\) is even, the remaining real factor at \(j=N/2\) is \((u_a+v_a)^N-1>0\).  Thus the total product is positive.
\end{proof}

\subsection{Power sums and exact small cases}

\begin{proposition}[power sums]\label{prop:powersums}
For every integer \(m\ge1\),
\begin{equation}\label{eq:powersum}
\sum_{j=0}^{N-1} Z_j^m
=N\sum_{r=0}^{m}
\binom{Nm}{Nr}(-1)^{N(m-r)}P^{m-r}Q^r.
\end{equation}
In particular,
\[
\Tr(Z)=N\bigl(Q+(-1)^NP\bigr),
\qquad \Phi_{N;P,Q}(0)=(P-Q)^N.
\]
\end{proposition}

\begin{proof}
Expand
\[
Z_j^m=P^m(\zeta_N^j\alpha-1)^{Nm}
=P^m\sum_{k=0}^{Nm}\binom{Nm}{k}(-1)^{Nm-k}\alpha^k\zeta_N^{jk}.
\]
Summing over \(j\) kills every term except \(k=Nr\).  Since \(\alpha^{Nr}=(Q/P)^r\), \eqref{eq:powersum} follows.  The constant term follows either from the same expansion or from
\(\prod_j(\zeta_N^jv-u)=(-1)^{N+1}(Q-P)\).
\end{proof}

Newton identities turn \eqref{eq:powersum} into a radical-free algorithm for every coefficient.  The first two nontrivial polynomials are
\begin{align*}
\Phi_{2;P,Q}(T)
 &=T^2-2(P+Q)T+(P-Q)^2,\\
\Phi_{3;P,Q}(T)
 &=T^3+3(P-Q)T^2+3(P^2+7PQ+Q^2)T+(P-Q)^3.
\end{align*}
These formulas provide useful unit tests for a Lean or computer-algebra implementation.

\section{The distinguished embedding: a scale-free first jet}

Define
\begin{equation}\label{eq:ca}
c_a=h_a'(1)=(a+1)\log(a+1)-a\log a>0.
\end{equation}
For our four pairs,
\begin{align*}
c_1&=2\log2,\\
c_2&=-2\log2+3\log3,\\
c_3&=8\log2-3\log3,\\
c_8&=-24\log2+18\log3.
\end{align*}

\begin{proposition}[uniform distinguished expansion]\label{prop:distinguished}
As \(\eps\to0\), uniformly for \(N\ge2\),
\begin{equation}\label{eq:dist-expansion}
\delta_a(N,q)
=c_a\eps+\frac{c_a^2}{2}\eps^2
+O_a\!\left(\frac{\eps^2}{N}+|\eps|^3\right).
\end{equation}
In particular, along lower convergents,
\(\delta_a(N,q)=c_a\eps(1+o(1))\).
\end{proposition}

\begin{proof}
Let \(L_a(t)=\log h_a(1+t)\).  Since \(h_a(1)=1\),
\[
L_a(t)=c_at+\frac{L_a''(0)}2t^2+O_a(t^3).
\]
With \(t=\eps/N\),
\[
N L_a(\eps/N)
=c_a\eps+O_a(\eps^2/N).
\]
Expanding \(\exp(NL_a)-1\) to second order gives \eqref{eq:dist-expansion}.
\end{proof}

The decisive feature is that the first-order size is \(c_a\eps\), not
\((a+1)^N\eps\).  The adjacent unit gap has removed the exponential scale from the distinguished conjugate.

\section{Exact factorization and the \texorpdfstring{$\zeta(3/2)$}{zeta(3/2)} boundary law}\label{sec:boundary}

\subsection{An exact archimedean factorization}
With \(z_j=\zeta_N^jv-u\), one has
\[
\prod_{j=0}^{N-1}z_j=(-1)^{N+1}(Q-P),
\]
so \(\prod_j z_j^N=(Q-P)^N\).  Pairing complex conjugates gives the positive exact identity
\begin{equation}\label{eq:factorization}
\mathscr R_N(P,Q)
=(Q-P)^N\prod_{j=0}^{N-1}\left|1-z_j^{-N}\right|.
\end{equation}
The factor \(j=0\) is asymptotic to \(c_a\eps\).  The remaining factors form a boundary layer around \(j=0\) and \(j=N\).

\subsection{The boundary-layer theorem}

\begin{theorem}[universal boundary law]\label{thm:boundary}
Let \(N/q<\beta\) run through lower convergents and fix \(a\in\mathcal A\).  Then
\begin{equation}\label{eq:boundary-law}
\log\prod_{j=1}^{N-1}\left|1-z_{a,j}^{-N}\right|
=-\gamma_a\sqrt N+o(\sqrt N),
\end{equation}
where
\begin{equation}\label{eq:gamma}
\gamma_a=\frac{\zR(3/2)}{\sqrt{2\pi a(a+1)}}.
\end{equation}
Consequently,
\begin{equation}\label{eq:full-asymptotic}
\boxed{\quad
\log\mathscr R_a(N,q)
=N\log(Q_a-P_a)+\log(c_a\eps)
-\gamma_a\sqrt N+o(\sqrt N).
\quad}
\end{equation}
\end{theorem}

\begin{proof}[Proof, with the domination details isolated]
Write \(b=a+1\) and \(\theta_j=2\pi j/N\).  Lower convergents give
\(s=1+O(N^{-2})\), hence \(u_a=a+O_a(N^{-2})\) and
\(v_a=b+O_a(N^{-2})\).  For \(j/\sqrt N\) bounded, the principal logarithm near \(1\) gives
\begin{equation}\label{eq:local-log}
N\log(v_ae^{i\theta_j}-u_a)
=2\pi i b j+2\pi^2ab\frac{j^2}{N}
+O_a\!\left(\frac{j^3}{N^2}+\frac1N\right).
\end{equation}
Because \(b\in\ZZ\), the first term contributes the phase \(1\) after exponentiation.  Thus, whenever \(j/\sqrt N\to x\ge0\),
\[
(v_ae^{i\theta_j}-u_a)^{-N}
\longrightarrow e^{-2\pi^2abx^2}.
\]
The same limit holds at the conjugate boundary \(N-j\).

For domination, the exact identity
\[
|v_ae^{i\theta}-u_a|^2
=(v_a-u_a)^2+4u_av_a\sin^2(\theta/2)
\]
shows, for \(1\le j\le N/2\),
\[
|z_{a,j}|^{-N}\le \exp(-c_a'j^2/N)
\]
with a fixed \(c_a'>0\) for all sufficiently large \(N\).  Since
\[
\log(1-|w|)\le\log|1-w|\le\log(1+|w|),
\]
the scaled summands are dominated by an integrable envelope comparable near zero to \(|\log x|\), and at infinity to \(e^{-c x^2}\).  Truncation followed by dominated Riemann-sum convergence therefore yields
\begin{align*}
\lim_{N\to\infty}\frac1{\sqrt N}
\sum_{j=1}^{N-1}\log|1-z_{a,j}^{-N}|
&=2\int_0^\infty\log(1-e^{-2\pi^2abx^2})\,dx\\
&=-\frac{\zR(3/2)}{\sqrt{2\pi ab}}.
\end{align*}
The integral is evaluated by expanding
\(\log(1-e^{-cx^2})=-\sum_{m\ge1}e^{-mcx^2}/m\) and integrating termwise.  Finally, the \(j=0\) factor in \eqref{eq:factorization} is
\[
1-h_a(s)^{-N}=c_a\eps(1+o(1))
\]
by \cref{prop:distinguished}, which proves \eqref{eq:full-asymptotic}.
\end{proof}

\begin{remark}[why the square root appears]
There are \(\asymp\sqrt N\) roots of unity within angular distance \(O(N^{-1/2})\) of \(1\).  For those roots, \(|z_j|^N\) is neither exponentially large nor close to a fixed constant; it varies on the Gaussian scale \(e^{2\pi^2a(a+1)j^2/N}\).  Their cumulative logarithm is therefore a Riemann sum with mesh \(N^{-1/2}\), and the Bose-type integral produces \(\zR(3/2)\).
\end{remark}

\subsection{Numerical check of the universal constant}
The following table uses the artificial but correctly scaled perturbation
\(s=1+N^{-2}\) and reports
\(-N^{-1/2}\log\prod_{j=1}^{N-1}|1-z_j^{-N}|\).  It is not used in the proof.

\begin{table}[ht]
\centering
\begin{tabular}{c@{\qquad}c@{\qquad}ccc}
\toprule
\(a\) & predicted \(\gamma_a\) & \(N=200\) & \(N=800\) & \(N=3200\)\\
\midrule
1 & 0.736937 & 0.356588 & 0.497187 & 0.592403 \\
2 & 0.425471 & 0.126603 & 0.226610 & 0.301414 \\
3 & 0.300853 & 0.052122 & 0.127170 & 0.189410 \\
8 & 0.122823 & 0.000051 & 0.013106 & 0.043425 \\
\bottomrule
\end{tabular}
\caption{Convergence to the boundary constant in \cref{thm:boundary}.}
\end{table}

Numerically,
\[
(\gamma_1,\gamma_2,\gamma_3,\gamma_8)
=(\GammaOne,\GammaTwo,\GammaThree,\GammaEight)\ldots.
\]

\section{The unique balanced four-Catalan quotient}\label{sec:balanced}

\subsection{The four integer gaps}
Put
\[
X=M^q=2^{N+\eps},\qquad Y=A^q=3^{N+\eps}.
\]
For the four pairs, the integer data \((P_a,Q_a)\) and the principal gaps are
\begin{center}
\begin{tabular}{ccl}
\toprule
\(a\) & \((P_a,Q_a)\) & \(Q_a-P_a\)\\
\midrule
1 & \((1,X)\) & \(X-1\)\\
2 & \((X,Y)\) & \(Y-X\)\\
3 & \((Y,X^2)\) & \(X^2-Y\)\\
8 & \((X^3,Y^2)\) & \(Y^2-X^3\)\\
\bottomrule
\end{tabular}
\end{center}
The endpoint heights are respectively
\(\log2,\log3,\log4,\log9\).

\subsection{Uniqueness of the balance vector}
Consider a multiplicative combination
\(\prod_{a\in\mathcal A}\mathscr R_a^{w_a}\) with integer exponents \(w_a\).  Cancellation of the two \(N^2\)-height components requires
\[
w_1+2w_3=0,\qquad w_2+2w_8=0.
\]
Each resultant contains one factor \(c_a\eps\), so cancellation of \(\log\eps\) requires
\(w_1+w_2+w_3+w_8=0\).  The primitive integer solution is unique up to sign:
\begin{equation}\label{eq:balance-vector}
(w_1,w_2,w_3,w_8)=(2,-2,-1,1).
\end{equation}
This yields the canonical quotient
\begin{equation}\label{eq:theta}
\Theta_{N,q}
=\frac{\mathscr R_1(N,q)^2\mathscr R_8(N,q)}
{\mathscr R_3(N,q)\mathscr R_2(N,q)^2}>0.
\end{equation}

\begin{theorem}[balanced quotient asymptotic]\label{thm:balanced}
Along lower convergents \(N/q<\beta\),
\begin{equation}\label{eq:theta-asymptotic}
\boxed{\quad
\Theta_{N,q}
=C_0\exp(-\kappa\sqrt N+o(\sqrt N)),
\quad}
\end{equation}
where
\begin{align}
C_0&=\frac{c_1^2c_8}{c_3c_2^2}=\BalancedConstant\ldots,\\
\kappa&=2\gamma_1+\gamma_8-\gamma_3-2\gamma_2\label{eq:kappa-def}\\
&=\frac{\zR(3/2)}{\sqrt{2\pi}}
\left(
\frac2{\sqrt2}+\frac1{\sqrt{72}}
-\frac1{\sqrt{12}}-\frac2{\sqrt6}
\right)
=\KappaBalanced\ldots>0.
\end{align}
In particular, \(0<\Theta_{N,q}<1\) for all sufficiently large lower convergents.
\end{theorem}

\begin{proof}
Insert \eqref{eq:full-asymptotic} with weights \eqref{eq:balance-vector}.  The \(\log\eps\) terms cancel and the constants give \(\log C_0\).  It remains only to verify that the exact gap terms cancel up to \(o(1)\).  Their weighted logarithm is
\begin{align*}
&2\log(X-1)+\log(Y^2-X^3)-\log(X^2-Y)-2\log(Y-X)\\
&=2\log(1-X^{-1})+\log(1-X^3/Y^2)
 -\log(1-Y/X^2)-2\log(1-X/Y).
\end{align*}
Here
\[
X^{-1}=O(2^{-N}),\quad Y/X^2=O((3/4)^N),\quad
X/Y=O((2/3)^N),\quad X^3/Y^2=O((8/9)^N).
\]
Multiplication by \(N\) still gives \(o(1)\).  The remaining square-root term is \(-\kappa\sqrt N\).  Positivity of \(\kappa\) follows directly from the displayed radical expression.
\end{proof}

Two simpler channels are also useful:
\begin{align}
\frac{\mathscr R_1^2}{\mathscr R_3}
&=\frac{c_1^2}{c_3}\eps
\exp\bigl(-(2\gamma_1-\gamma_3)\sqrt N+o(\sqrt N)\bigr),\label{eq:channel13}\\
\frac{\mathscr R_2^2}{\mathscr R_8}
&=\frac{c_2^2}{c_8}\eps
\exp\bigl(-(2\gamma_2-\gamma_8)\sqrt N+o(\sqrt N)\bigr),\label{eq:channel28}
\end{align}
with constants
\(2\gamma_1-\gamma_3=\KappaThirteen\ldots\) and
\(2\gamma_2-\gamma_8=\KappaTwoEight\ldots\).  Dividing \eqref{eq:channel13} by \eqref{eq:channel28} eliminates \(\eps\) and recovers \cref{thm:balanced}.

\subsection{The exact primitive-denominator obstruction}
Write \(\Theta_{N,q}=U_{N,q}/V_{N,q}\) in lowest terms.  Explicitly,
\begin{equation}\label{eq:reduced-denom}
V_{N,q}
=\frac{\mathscr R_3\mathscr R_2^2}
{\gcd(\mathscr R_1^2\mathscr R_8,\,\mathscr R_3\mathscr R_2^2)}.
\end{equation}

\begin{corollary}[necessary denominator explosion]\label{cor:denominator}
Every hypothetical nonintegral solution forces
\begin{equation}\label{eq:denom-liminf}
\liminf_{\substack{N/q<\beta\\\text{lower convergents}}}
\frac{\log V_{N,q}}{\sqrt N}\ge\kappa.
\end{equation}
\end{corollary}

\begin{proof}
The reduced numerator \(U_{N,q}\) is a positive integer, so
\(V_{N,q}\ge1/\Theta_{N,q}\).  Apply \cref{thm:balanced}.
\end{proof}

\begin{criterion}[cross-resultant gcd criterion]\label{crit:gcd}
The Alaoglu--Erd\H{o}s conjecture follows if one proves that, for every hypothetical primitive pair \((M,A)\), there are \(\eta>0\) and infinitely many lower convergents \(N/q\) for which
\begin{equation}\label{eq:gcd-criterion}
\frac{\mathscr R_3\mathscr R_2^2}
{\gcd(\mathscr R_1^2\mathscr R_8,\,\mathscr R_3\mathscr R_2^2)}
\le \exp((\kappa-\eta)\sqrt N).
\end{equation}
In particular, divisibility
\(\mathscr R_3\mathscr R_2^2\mid\mathscr R_1^2\mathscr R_8\)
for one sufficiently large lower convergent would finish the proof.
\end{criterion}

This is a severe divisibility target: the individual resultants have logarithms of order \(N^2\), whereas the permitted primitive denominator is only \(e^{O(\sqrt N)}\).  The point is not that the target is easy; it is that the exact amount of arithmetic cancellation needed is now known, with a universal constant.

\section{The Catalan tangent lattice}\label{sec:tangent}

\subsection{Exact real identities}
Set
\[
U=2^s,\qquad V=3^s
\]
and define the unpowered defects
\begin{align*}
e_1&=U-2,\\
e_2&=V-U-1,\\
e_3&=U^2-V-1,\\
e_8&=V^2-U^3-1.
\end{align*}
Writing \(E=e_1\), \(F=e_2\), so that \(U=2+E\), \(V=3+E+F\), gives the exact identities
\begin{align}
-3e_1+e_2+e_3&=e_1^2,\label{eq:exact-tangent1}\\
6e_1-6e_2+e_8&=-e_1^3-5e_1^2+2e_1e_2+e_2^2.\label{eq:exact-tangent2}
\end{align}
Thus two integer combinations kill the entire first jet at \(s=1\), not merely a numerically fitted coefficient.

\begin{proposition}[the full integer tangent lattice]\label{prop:tangent-lattice}
The lattice of integer relations among \(c_1,c_2,c_3,c_8\) is
\begin{equation}\label{eq:tangent-lattice}
\left\{r\in\ZZ^4:\sum_{a\in\mathcal A}r_ac_a=0\right\}
=\ZZ(-3,1,1,0)\oplus\ZZ(6,-6,0,1).
\end{equation}
\end{proposition}

\begin{proof}
Separate coefficients of \(\log2\) and \(\log3\), which are rationally independent.  The two equations are
\[
2r_1-2r_2+8r_3-24r_8=0,
\qquad 3r_2-3r_3+18r_8=0.
\]
Writing \(r_3=t\), \(r_8=u\) gives
\[
(r_1,r_2,r_3,r_8)=t(-3,1,1,0)+u(6,-6,0,1).
\]
\end{proof}

For the powered distinguished defects \(\delta_a=(1+e_a)^N-1\), \cref{prop:distinguished} gives
\begin{align}
-3\delta_1+\delta_2+\delta_3&=O_\beta(\eps^2),\label{eq:powered-cancel1}\\
6\delta_1-6\delta_2+\delta_8&=O_\beta(\eps^2).
\label{eq:powered-cancel2}
\end{align}
More explicitly, the quadratic coefficients are
\begin{align*}
\frac12(-3c_1^2+c_2^2+c_3^2)
&=28(\log2)^2-30\log2\log3+9(\log3)^2,\\
\frac12(6c_1^2-6c_2^2+c_8^2)
&=9\bigl(32(\log2)^2-44\log2\log3+15(\log3)^2\bigr).
\end{align*}

\subsection{The Kummer-rank tax}
In the full Kummer field generated by \(U,V\), the four descended defects use the monomial directions
\begin{equation}\label{eq:kummer-vectors}
\nu_1=(1,0),\quad
\nu_2=(-1,1),\quad
\nu_3=(2,-1),\quad
\nu_8=(-3,2).
\end{equation}
They correspond respectively to \(U\), \(V/U\), \(U^2/V\), and \(V^2/U^3\).

\begin{proposition}[cancellation necessarily restores rank two]\label{prop:rank-tax}
Every nonzero integer tangent relation in \eqref{eq:tangent-lattice} has support containing at least three of the four directions in \eqref{eq:kummer-vectors}, and those directions span \(\ZZ^2\).  Hence a first-jet-cancelled additive certificate cannot remain in a single generic degree-\(N\) Kummer subfield; absent an additional arithmetic collapse, it re-enters a rank-two extension of degree at most \(N^2\).
\end{proposition}

\begin{proof}
The parameterization in \cref{prop:tangent-lattice} shows that no nonzero relation has support of size one or two.  No two vectors in \eqref{eq:kummer-vectors} are collinear, so any support of size at least three spans rank two.
\end{proof}

There is a second incompatibility.  Associate endpoint-height vectors
\[
\eta_1=(1,0),\quad\eta_2=(0,1),\quad
\eta_3=(2,0),\quad\eta_8=(0,2),
\]
representing \(\log2,\log3,\log4,\log9\).  If a tangent relation
\(r=t(-3,1,1,0)+u(6,-6,0,1)\) also cancelled endpoint height, then
\[
r_1+2r_3=-t+6u=0,
\qquad r_2+2r_8=t-4u=0,
\]
so \(t=u=0\).  Thus no nonzero same-weight linear certificate cancels both the first real jet and the leading two-dimensional height.  The balanced quotient of \cref{sec:balanced} avoids this obstruction by using multiplicative norm weights rather than an additive tangent relation.

\section{Prime-degree Frobenius collapse and the finite logarithm}\label{sec:frobenius}

\subsection{Universal collapse modulo the prime degree}

\begin{theorem}[Frobenius collapse]\label{thm:frob-collapse}
Let \(\ell\) be prime.  As a polynomial identity over \(\FF_\ell[P,Q,T]\),
\begin{equation}\label{eq:frob-poly}
\Phi_{\ell;P,Q}(T)\equiv(T-Q+P)^\ell\pmod\ell.
\end{equation}
Consequently,
\[
\Phi_{\ell;P,Q}(1)\equiv(1-Q+P)^\ell\pmod\ell.
\]
\end{theorem}

\begin{proof}
In characteristic \(\ell\), \((X-1)^\ell=X^\ell-1\).  On the scheme cut out by \(PX^\ell-Q\), the second resultant polynomial
\(T-P(X-1)^\ell\) becomes the constant \(T-Q+P\), with total multiplicity \(\ell\).  The identity also follows coefficientwise from \cref{prop:powersums} and Newton identities.
\end{proof}

The congruence \(Q-P\equiv1\pmod\ell\) is the finite-field shadow of the exact adjacent relation at \(s=1\).  It yields much more than a single factor of \(\ell\).

\begin{theorem}[exact first Frobenius quotient]\label{thm:first-quotient}
Let \(\ell\) be an odd prime, \(\ell\nmid P\), and assume
\(Q-P-1=\ell C\).  Define
\begin{equation}\label{eq:Hpoly}
H_\ell(X)=\frac{(X-1)^\ell-X^\ell+1}{\ell}\in\ZZ[X].
\end{equation}
Then
\begin{equation}\label{eq:frob-factorization}
\mathscr R_\ell(P,Q)
=P^\ell\ell^\ell
\Res_X\!\left(
X^\ell-\frac QP,\;
H_\ell(X)+\frac CP
\right)
\end{equation}
inside \(\ZZ_\ell\).  In particular,
\begin{equation}\label{eq:ell-ell-div}
\ell^\ell\mid\mathscr R_\ell(P,Q).
\end{equation}
Moreover, with \(r=Q/P\in\FF_\ell\),
\begin{equation}\label{eq:first-residue}
\frac{\mathscr R_\ell(P,Q)}{\ell^\ell}
\equiv P H_\ell(r)+C\pmod\ell.
\end{equation}
For odd \(\ell\),
\begin{equation}\label{eq:finite-log}
H_\ell(X)\equiv\sum_{k=1}^{\ell-1}\frac{X^k}{k}\pmod\ell.
\end{equation}
\end{theorem}

\begin{proof}
Over \(\ZZ_\ell\), put
\[
f=X^\ell-Q/P,\qquad g=(X-1)^\ell-1/P.
\]
The hypothesis gives
\[
g-f=\ell\left(H_\ell(X)+C/P\right).
\]
Adding a multiple of the monic polynomial \(f\) does not change the resultant, and scaling the second polynomial by \(\ell\) multiplies the resultant by \(\ell^\ell\).  Rescaling from \(f,g\) to the original polynomials contributes \(P^\ell\), yielding \eqref{eq:frob-factorization}.  Modulo \(\ell\),
\(f=(X-r)^\ell\), so the remaining resultant is the \(\ell\)-th power of the second polynomial evaluated at \(r\), which equals that value in \(\FF_\ell\).  Finally,
\(\ell^{-1}\binom\ell k\equiv(-1)^{k-1}/k\pmod\ell\); the signs in \eqref{eq:Hpoly} cancel and give \eqref{eq:finite-log}.
\end{proof}

For example, at \(P=1,Q=2\),
\[
H_\ell(2)=\frac{2-2^\ell}{\ell}
\equiv-2q_\ell(2)\pmod\ell,
\qquad q_\ell(2)=\frac{2^{\ell-1}-1}{\ell}.
\]
Thus the generic valuation is exactly \(\ell\); an extra factor is a Wieferich-type event.

\subsection{Simultaneous resonance and the adelic tangent lattice}
Now specialize to a prime lower-convergent numerator \(N=\ell>3\).  Suppose the strongest natural first-order resonance holds:
\begin{equation}\label{eq:sim-resonance}
X=M^q\equiv2\pmod\ell,
\qquad Y=A^q\equiv3\pmod\ell.
\end{equation}
Write modulo \(\ell^2\)
\[
X=2+\ell x,
\qquad Y=3+\ell y,
\]
and let
\[
q_\ell(r)=\frac{r^{\ell-1}-1}{\ell}\pmod\ell
\quad(r=2,3)
\]
be Fermat quotients.  Since all four adjacent congruences then hold, define
\[
\widehat{\mathscr R}_a=\mathscr R_a/\ell^\ell\in\ZZ.
\]

\begin{theorem}[four first-quotient residues]\label{thm:Wtable}
Under \eqref{eq:sim-resonance}, the residues
\(W_a=\widehat{\mathscr R}_a\pmod\ell\) are
\begin{align}
W_1&=-2q_\ell(2)+x,\label{eq:W1}\\
W_2&=2q_\ell(2)-3q_\ell(3)+y-x,\label{eq:W2}\\
W_3&=-8q_\ell(2)+3q_\ell(3)+4x-y,\label{eq:W3}\\
W_8&=24q_\ell(2)-18q_\ell(3)+6y-12x.
\label{eq:W8}
\end{align}
They satisfy the exact tangent relations
\begin{equation}\label{eq:Wrelations}
-3W_1+W_2+W_3=0,
\qquad
6W_1-6W_2+W_8=0
\quad\text{in }\FF_\ell.
\end{equation}
\end{theorem}

\begin{proof}
Apply \eqref{eq:first-residue} to the four pairs
\[
(1,X),\quad(X,Y),\quad(Y,X^2),\quad(X^3,Y^2).
\]
For \(a,b,a-b\not\equiv0\pmod\ell\), direct expansion gives
\begin{equation}\label{eq:H-rational}
H_\ell(a/b)
\equiv\frac{(a-b)q_\ell(a-b)-a q_\ell(a)+bq_\ell(b)}{b}
\pmod\ell.
\end{equation}
Use \(q_\ell(4)=2q_\ell(2)\),
\(q_\ell(8)=3q_\ell(2)\), and
\(q_\ell(9)=2q_\ell(3)\).  The four values are
\begin{align*}
H_\ell(2)&=-2q_\ell(2),\\
2H_\ell(3/2)&=2q_\ell(2)-3q_\ell(3),\\
3H_\ell(4/3)&=-8q_\ell(2)+3q_\ell(3),\\
8H_\ell(9/8)&=24q_\ell(2)-18q_\ell(3).
\end{align*}
The first-order corrections \((Q-P-1)/\ell\) are respectively
\(x\), \(y-x\), \(4x-y\), and \(6y-12x\).  This proves the table; substitution proves \eqref{eq:Wrelations}.
\end{proof}

The same two lattice vectors therefore govern
\begin{itemize}
\item the real derivatives \((c_1,c_2,c_3,c_8)\), with \(\log2,\log3\), and
\item the first finite-place quotients \((W_1,W_2,W_3,W_8)\), with Fermat quotients and residue perturbations.
\end{itemize}
This is an exact adelic parallel, not a metaphor.

\begin{corollary}[one-variable residual map]\label{cor:belyi}
Assume additionally that \(W_1W_2W_3W_8\ne0\), and set
\(t=W_2/W_1\in\FF_\ell\).  The common factors \(\ell^{3\ell}\) cancel between numerator and denominator of the balanced quotient, and
\begin{equation}\label{eq:belyi}
\Theta_{\ell,q}
\equiv\frac{6(t-1)}{t^2(3-t)}\pmod\ell.
\end{equation}
\end{corollary}

\begin{proof}
From \eqref{eq:Wrelations},
\(W_3=W_1(3-t)\) and \(W_8=6W_1(t-1)\).  Substitute into
\(W_1^2W_8/(W_3W_2^2)\).
\end{proof}

The exceptional projective values \(t=0,1,3,\infty\) mark precisely where one of the first quotients acquires an extra factor of \(\ell\).  Thus the local obstruction has been reduced to the recurrence and distribution of a concrete finite-logarithm parameter.

\subsection{Quantitative audit of the Frobenius gain}
For one resultant, \cref{thm:boundary} gives
\[
\log\mathscr R_a\sim N^2\log(a+1),
\]
whereas the first Frobenius rank drop supplies only
\(\log(\ell^\ell)=\ell\log\ell=o(\ell^2)\).  Even four simultaneous first-layer collapses cannot defeat an individual norm's main height.  In the balanced quotient, however, both the archimedean \(N^2\) terms and the common \(\ell^{3\ell}\) finite-place terms cancel.  The meaningful arithmetic begins at the first quotients \(W_a\), exactly where \cref{thm:Wtable,cor:belyi} place it.

\section{What has genuinely changed, and what has not}\label{sec:audit}

\subsection{New results relative to the unified report}
The following items are not restatements of the inherited Kummer/product-formula analyses.
\begin{enumerate}
\item \textbf{Diagonal, rather than generic, Kummer descent.}
The powered adjacent difference lies in \(\QQ((Q/P)^{1/N})\), degree at most \(N\), even though the two radicals separately can require degree \(N^2\).

\item \textbf{A universal integral polynomial.}
The exact polynomial \(\Phi_{N;P,Q}\), its sparse resultant, and its closed power sums make the construction computable without choosing radicals.

\item \textbf{A sharp conjugate audit.}
The other embeddings are not merely bounded by a house.  Their exact contribution has the subleading law
\(-\zR(3/2)\sqrt N/\sqrt{2\pi a(a+1)}\).

\item \textbf{The unique balanced quotient.}
The four Catalan resultants have one primitive multiplicative combination that cancels the full two-dimensional endpoint height and the unknown approximation error.  It leaves a universal stretched-exponential decay.

\item \textbf{A precise denominator threshold.}
A counterexample forces \(V_{N,q}\ge e^{(\kappa-o(1))\sqrt N}\).  Any fixed improvement in the opposite direction proves the conjecture.

\item \textbf{An exact local first quotient.}
The \(\ell^\ell\) Frobenius factor and the truncated-logarithm residue are explicit, and the four residues reproduce the real tangent lattice exactly.

\item \textbf{A one-parameter finite-field reduction.}
After common first-layer cancellation, the balanced quotient is governed by
\(6(t-1)/(t^2(3-t))\).
\end{enumerate}

\subsection{The remaining bottleneck}
The balanced quotient is rational, not known to be integral.  Its denominator can in principle be enormous because the four resultants are values of different specializations of the universal Catalan resultant.  The exact real coupling through \(X=M^q\), \(Y=A^q\) has not yet been converted into the nearly complete common divisibility demanded by \eqref{eq:gcd-criterion}.

This is now the central question:

\begin{problem}[cross-Catalan primitive divisor problem]\label{prob:cross}
For
\[
R_1=\mathscr R_N(1,X),\quad
R_2=\mathscr R_N(X,Y),\quad
R_3=\mathscr R_N(Y,X^2),\quad
R_8=\mathscr R_N(X^3,Y^2),
\]
where \(X=2^{N+\eps}\), \(Y=3^{N+\eps}\) are integers arising from a common exponent, control
\[
\frac{R_3R_2^2}{\gcd(R_1^2R_8,R_3R_2^2)}
\]
at the \(e^{c\sqrt N}\) scale.
\end{problem}

A generic polynomial-value gcd estimate is unlikely to suffice: genericity predicts many primitive prime factors, while a proof needs substantial cross-specialization sharing.  The common Kummer origin and the local tangent identities are the non-generic structure that must be exploited.

\section{Concrete next attacks}\label{sec:next}

\subsection{Prime-ideal transfer in the full rank-two Kummer field}
Let \(K=\QQ(\zeta_N,U,V)\), where \(U^N=X\), \(V^N=Y\), and define
\begin{align*}
D_1&=(U-1)^N-1,\\
D_2&=(V-U)^N-1,\\
D_3&=(U^2-V)^N-1,\\
D_8&=(V^2-U^3)^N-1.
\end{align*}
The full norm of the rational function
\begin{equation}\label{eq:alpha-full}
\mathcal U=\frac{D_1^2D_8}{D_3D_2^2}
\end{equation}
contains \(\Theta_{N,q}^N\), up to the predictable multiplicities of the one-dimensional quotient fields.  A proof that the denominator ideal of \(\mathcal U\) has norm
\(\exp(o(N^{3/2}))\) would contradict \cref{thm:balanced}.  The exact identities \eqref{eq:exact-tangent1}--\eqref{eq:exact-tangent2} suggest studying prime ideals where one or more \(D_a\) vanish, separated into residue characteristic dividing \(N\) and prime-to-\(N\) cases.

\subsection{Finite-log dynamics at prime numerators}
Under \eqref{eq:sim-resonance}, all first-layer information is encoded by
\[
t=\frac{2q_\ell(2)-3q_\ell(3)+y-x}{-2q_\ell(2)+x}.
\]
The exceptional values \(0,1,3,\infty\) control extra numerator or denominator valuation.  A useful theorem would show that repeated denominator-favoring exceptional values force a multiplicative dependence among \(2,3,M,A\), or else occur too sparsely to create the denominator growth required by \cref{cor:denominator}.  This is a finite-field problem with explicit inputs, not an unspecified appeal to ``random residues.''

\subsection{Subresultant identities among the four specializations}
The universal curve
\[
\mathscr R_N(P,Q)=0
\]
is the implicit equation of \(P=t^N\), \(Q=(t+1)^N\).  The four values in \cref{prob:cross} are evaluations along a chain of monomial maps.  Computing multivariate subresultants and syzygies for small \(N\) may reveal a stable identity whose content controls the cross gcd.  The appropriate search object is not an equality of the four resultants, but an ideal containment after saturating by \(XY(X-Y)(X^2-Y)(Y^2-X^3)\).

\subsection{Boundary law with arithmetic weights}
The proof of \cref{thm:boundary} treats all roots of unity equally.  A weighted version, retaining residue classes or prime-ideal decomposition types, could connect the \(\sqrt N\) archimedean boundary directly to the finite-log parameter in \cref{cor:belyi}.  The desired statement would be an adelic local limit theorem for the same narrow arc around \(1\).

\subsection{A realistic computational campaign}
For moderate \(N\), compute \(\mathscr R_N(P,Q)\) by modular sparse resultants and Chinese remaindering, not by expanding radicals.  For synthetic integer pairs \((X,Y)\) close to \((2^N,3^N)\), record
\begin{itemize}
\item factor overlap among \(R_1^2R_8\) and \(R_3R_2^2\),
\item the reduced denominator on the \(\sqrt N\) scale,
\item prime-degree first quotients and the parameter \(t\), and
\item whether exceptional local residues correlate with unusually large cross gcd.
\end{itemize}
Synthetic pairs cannot prove the conjecture, but they can falsify overoptimistic universal divisibility claims quickly and identify the correct saturation factors for a symbolic theorem.

\section{Lean formalization blueprint}\label{sec:lean}

The algebraic core is substantially easier to formalize than the original transcendence problem.  A useful division into modules is as follows.

\subsection{Module A: universal resultants}
\begin{enumerate}
\item Define the sparse polynomials
\[
f_{N,P,Q}(X)=PX^N-Q,
\qquad g_{N,P,T}(X)=T-P(X-1)^N.
\]
\item Define \(\Phi\) through the resultant and prove coefficientwise divisibility by \(P^N\), or equivalently construct it through the Galois-stable multiset of integral roots and then prove \eqref{eq:resultant}.
\item Prove \cref{prop:powersums} using the roots-of-unity filter in \(\mathrm{ZMod}\,N\), then derive the first coefficients by Newton identities.
\item Define \(\mathscr R_N(P,Q)=(-1)^N\Phi(1)\) and prove positivity for the real adjacent-pair specialization.
\end{enumerate}

Indicative theorem statements are
\begin{lstlisting}[language={},basicstyle=\ttfamily\small,frame=single,breaklines=true]
theorem cyclicPoly_mem_int
  (N P Q : Nat) (hN : 0 < N) (hP : 0 < P) :
  cyclicPoly N P Q \in Polynomial.Z := ...

theorem cyclic_powerSum
  (m : Nat) :
  powerSum (cyclicRoots N P Q) m =
    N * \sum r in Finset.range (m+1),
      choose (N*m) (N*r) * (-1)^(N*(m-r)) * P^(m-r) * Q^r := ...
\end{lstlisting}

\subsection{Module B: real distinguished identities}
Formalize \eqref{eq:exact-tangent1}--\eqref{eq:exact-tangent2} as polynomial identities in two variables before introducing exponentials.  Prove the tangent lattice by elementary integer linear algebra.  The expansion \eqref{eq:dist-expansion} requires only Taylor's theorem for \(\exp\) and \(\log\) on a compact neighborhood of \(1\).

\subsection{Module C: Frobenius and finite logarithms}
\begin{enumerate}
\item Prove \((X-1)^\ell=X^\ell-1\) in \(\FF_\ell[X]\).
\item Derive \eqref{eq:frob-poly} through resultant functoriality under coefficient reduction.
\item Define \(H_\ell\) over \(\ZZ[X]\) using divisibility of middle binomial coefficients by \(\ell\).
\item Prove resultant invariance under adding a multiple of the first polynomial and the scaling law needed for \eqref{eq:frob-factorization}.
\item Formalize the rational evaluation formula \eqref{eq:H-rational}, then the four-row table and its two linear relations by \texttt{ring}.
\end{enumerate}
These are excellent early targets for a rival-proof race because they are exact, finite, and do not depend on delicate analysis.

\subsection{Module D: the boundary law}
The analytic theorem is the longest formalization component.  A staged approach is preferable.
\begin{enumerate}
\item First prove two-sided bounds
\[
e^{-(\gamma_a+o(1))\sqrt N}
\le\prod_{j=1}^{N-1}|1-z_{a,j}^{-N}|
\le e^{-(\gamma_a-o(1))\sqrt N}
\]
using a truncated window \(j\le L\sqrt N\), an explicit Gaussian tail, and then \(L\to\infty\).
\item Isolate the integrable singularity of \(\log(1-e^{-cx^2})\) by subtracting \(2\log x\) near zero.
\item Evaluate the integral through monotone convergence of the nonnegative series
\(-\log(1-e^{-cx^2})\).
\end{enumerate}
The final denominator criterion then follows from elementary properties of reduced fractions and \(\gcd\).

\section{Reference implementation sketch}

The following short SymPy program constructs \(\Phi_{N;P,Q}\) exactly for small \(N\), verifies the resultant normalization, and evaluates the cyclic resultant.  Large computations should replace symbolic expansion by modular subresultants.

\begin{lstlisting}[language=Python,basicstyle=\ttfamily\small,frame=single,breaklines=true,showstringspaces=false]
from sympy import symbols, resultant, expand, cancel

X, T, P, Q = symbols("X T P Q")

def cyclic_poly(n: int):
    f = P*X**n - Q
    g = T - P*(X - 1)**n
    r = resultant(f, g, X)
    phi = cancel(r / P**n)
    assert expand(phi).as_poly(T).LC() == 1
    return expand(phi)

def cyclic_resultant(n: int, p: int, q: int) -> int:
    phi = cyclic_poly(n).subs({P: p, Q: q, T: 1})
    return int((-1)**n * phi)

for n in (2, 3, 5):
    print(n, cyclic_poly(n))
\end{lstlisting}

For the boundary constant, evaluate logarithms rather than the enormous integer:
\begin{lstlisting}[language=Python,basicstyle=\ttfamily\small,frame=single,breaklines=true,showstringspaces=false]
import cmath, math

def boundary_log(a: int, n: int) -> float:
    s = 1.0 + 1.0/(n*n)
    u, v = a**s, (a+1)**s
    total = 0.0
    for j in range(1, n):
        theta = 2*math.pi*j/n
        z = v*cmath.exp(1j*theta) - u
        total += math.log(abs(1 - cmath.exp(-n*cmath.log(z))))
    return total
\end{lstlisting}

\section{Conclusion}

The adjacent-smooth identities do more than provide four small numbers.  They organize into one coherent structure:
\begin{enumerate}
\item each unit gap gives a scale-free distinguished defect;
\item powering by the convergent numerator kills the diagonal Kummer direction;
\item the full conjugate product is an explicit integer resultant;
\item the narrow arc of roots near \(1\) produces a universal \(\zR(3/2)\sqrt N\) boundary term;
\item the identities \(4=2^2\) and \(9=3^2\) select a unique balanced quotient in which the \(N^2\) height and the approximation error both disappear;
\item at prime degree, Frobenius collapse and the finite logarithm reproduce the same tangent lattice over \(\FF_\ell\).
\end{enumerate}

The resulting obstruction is exact:
\[
\Theta_{N,q}=C_0e^{-\kappa\sqrt N+o(\sqrt N)}
\]
is a positive rational built from four ordinary integers.  A counterexample can survive only by creating a primitive denominator of at least
\(e^{(\kappa-o(1))\sqrt N}\) after enormous cross-resultant cancellation.  The next proof attempt should therefore concentrate on prime-ideal transfer among the four specializations and on the finite-logarithm parameter \(t\), rather than returning to undifferentiated house bounds.

No proof of the conjecture is claimed here.  What has been achieved is a new, sharply quantified bridge between the real continued-fraction approximation, one-dimensional Kummer descent, cyclic-resultant asymptotics, and finite-place Frobenius geometry.  It is narrow enough to formalize and concrete enough to attack computationally.

\appendix

\section{A more explicit domination lemma for the boundary law}

For completeness, this appendix records a standard way to turn the proof of \cref{thm:boundary} into epsilon--delta estimates.

\begin{lemma}[Gaussian modulus bound]
For each fixed \(a\in\mathcal A\), there are constants \(c,C>0\) such that, for all sufficiently large lower-convergent numerators \(N\) and \(1\le j\le N/2\),
\[
\exp\!\left(c\frac{j^2}{N}\right)
\le |z_{a,j}|^N
\le \exp\!\left(C\left(\frac{j^2}{N}+\frac{j^4}{N^3}+\frac1N\right)\right)
\]
whenever \(j\le N^{3/4}\), while \(|z_{a,j}|^N\ge e^{cN^{1/2}}\) for \(j\ge N^{3/4}\).
\end{lemma}

\begin{proof}
Use
\[
|z_{a,j}|^2=(v_a-u_a)^2+4u_av_a\sin^2(\pi j/N)
\]
and \(v_a-u_a=1+O(N^{-2})\), \(u_av_a=a(a+1)+O(N^{-2})\).  On \([0,\pi/2]\), \(2x/\pi\le\sin x\le x\).  Taking \(N/2\) times the logarithm gives the stated estimates.
\end{proof}

Let
\[
F_N(j)=\log|1-z_{a,j}^{-N}|,
\qquad f_a(x)=\log(1-e^{-2\pi^2a(a+1)x^2}).
\]
For every fixed \(L\), \eqref{eq:local-log} gives uniform convergence
\[
F_N(j)=f_a(j/\sqrt N)+o(1)
\]
for \(1\le j\le L\sqrt N\), after pairing \(j\) with \(N-j\).  The Gaussian modulus bound gives
\[
|F_N(j)|\le
-\log(1-e^{-c j^2/N})+e^{-c j^2/N},
\]
whose scaled sum is bounded by an integrable function.  The tail
\(j>L\sqrt N\) is at most
\(C\sqrt N\int_L^\infty e^{-cx^2}\,dx\).  First let \(N\to\infty\), then \(L\to\infty\), obtaining \eqref{eq:boundary-law}.

\section{Universal coefficient facts}

The power sums in \eqref{eq:powersum} imply several useful structural properties.

\begin{proposition}
For prime \(\ell\), every non-leading, nonconstant coefficient of
\(\Phi_{\ell;P,Q}\) is divisible by \(\ell\), and the constant term is \((P-Q)^\ell\).  Equivalently,
\[
\Phi_{\ell;P,Q}(T)
=T^\ell+(P-Q)^\ell+\ell\,K_{\ell;P,Q}(T)
\]
for some \(K_{\ell;P,Q}\in\ZZ[T]\) of degree at most \(\ell-1\) and zero constant term.
\end{proposition}

\begin{proof}
This is another form of \cref{thm:frob-collapse}.  It can also be proved directly from Newton identities because all power sums of order \(<\ell\) are multiples of \(\ell\).
\end{proof}

For \(N=2\), the one-variable specialization already shows the geometric meaning:
\[
\mathscr R_2(1,X)=X(X-4).
\]
The factor \(X-2^N\) persists for every \(N\), because
\(\mathscr R_N(1,2^N)=0\).  More generally, \(\mathscr R_N(P,Q)=0\) is the implicit Catalan curve obtained by eliminating \(t\) from
\(P=t^N\), \(Q=(t+1)^N\).  The four resultants in \cref{sec:balanced} are four monomial specializations of this same universal curve.

\section{Formal statement inventory}

For a proof assistant or independent referee, the new claims can be checked in the following order.

\begin{longtable}{p{0.08\textwidth}p{0.55\textwidth}p{0.27\textwidth}}
\toprule
ID & Statement & Dependency\\
\midrule
A1 & \(\Phi_{N;P,Q}\in\ZZ[T]\) and resultant formula \eqref{eq:resultant}. & Algebraic integers/resultants\\
A2 & Power sums \eqref{eq:powersum}; examples for \(N=2,3\). & Roots-of-unity filter\\
A3 & Positivity of \(\mathscr R_a\) for lower approximants. & Complex conjugation\\
R1 & Distinguished expansion \eqref{eq:dist-expansion}. & Taylor theorem\\
R2 & Boundary law \eqref{eq:boundary-law}. & Local log expansion, dominated Riemann sums\\
R3 & Balanced quotient asymptotic \eqref{eq:theta-asymptotic}. & R2 plus exact gap cancellation\\
R4 & Denominator lower bound \eqref{eq:denom-liminf}. & Reduced fractions\\
T1 & Exact real identities \eqref{eq:exact-tangent1}--\eqref{eq:exact-tangent2}. & Polynomial expansion\\
T2 & Tangent lattice \eqref{eq:tangent-lattice}; rank tax. & Integer linear algebra\\
F1 & Frobenius collapse \eqref{eq:frob-poly}. & Freshman's dream/resultants\\
F2 & \(\ell^\ell\) factor and residue \eqref{eq:first-residue}. & Resultant scaling\\
F3 & Four-residue table \eqref{eq:W1}--\eqref{eq:W8}. & Fermat quotients\\
F4 & Residual map \eqref{eq:belyi}. & Two linear identities\\
\bottomrule
\end{longtable}

\begin{thebibliography}{99}

\bibitem{Unified}
V. Reshetnikov and collaborators,
\emph{Unified progress report on the Alaoglu--Erd\H{o}s conjecture},
unpublished working report supplied with the prompt, 2026.

\bibitem{AlaogluErdos}
L. Alaoglu and P. Erd\H{o}s,
On highly composite and similar numbers,
\emph{Transactions of the American Mathematical Society} \textbf{56} (1944), 448--469.

\bibitem{Mihailescu}
P. Mih\u{a}ilescu,
Primary cyclotomic units and a proof of Catalan's conjecture,
\emph{Journal f\"ur die reine und angewandte Mathematik} \textbf{572} (2004), 167--195.

\bibitem{Matveev}
E. M. Matveev,
An explicit lower bound for a homogeneous rational linear form in logarithms of algebraic numbers. II,
\emph{Izvestiya: Mathematics} \textbf{64} (2000), 1217--1269.

\bibitem{Waldschmidt}
M. Waldschmidt,
\emph{Diophantine Approximation on Linear Algebraic Groups},
Springer, 2000.

\bibitem{Lang}
S. Lang,
\emph{Introduction to Transcendental Numbers},
Addison--Wesley, 1966.

\bibitem{Hillar}
C. J. Hillar,
Cyclic resultants,
\emph{Journal of Symbolic Computation} \textbf{39} (2005), 653--669.

\bibitem{GKZ}
I. M. Gel'fand, M. M. Kapranov, and A. V. Zelevinsky,
\emph{Discriminants, Resultants, and Multidimensional Determinants},
Birkh\"auser, 1994.

\bibitem{Washington}
L. C. Washington,
\emph{Introduction to Cyclotomic Fields}, 2nd ed.,
Springer, 1997.

\bibitem{ArxivSearch}
ArXiv,
search entry point for the reported 2026 Jacobian-conjecture manuscript,
\url{https://arxiv.org/search/?query=The+Jacobian+Conjecture+is+false&searchtype=title},
accessed 22 August 2026.

\bibitem{AnthropicResearch}
Anthropic,
public research index,
\url{https://www.anthropic.com/research},
accessed 22 August 2026.

\bibitem{ICARM30}
NSF Institute for Computer-Aided Reasoning in Mathematics,
Elliptic Curve Rank Leaderboard, curve~\#273,
\url{https://elliptic-rank.icarm.cloud/curve/273},
accessed 22 August 2026.

\bibitem{HadamardRepo}
\texttt{bbeartheancient/hoa64},
Hadamard Matrix Construction and Search Toolchain,
\url{https://github.com/bbeartheancient/hoa64},
accessed 22 August 2026.

\bibitem{Fable668}
By Claude,
``Seventy-Four Is More Than Fifty-Five,''
\url{https://byclaude.net/seventy-four-is-more-than-fifty-five},
24 July 2026; accessed 22 August 2026.

\bibitem{JacobianMoment}
C. D. Long,
Small counterexamples to the Gaussian moments conjecture,
\url{https://arxiv.org/abs/2607.18186},
accessed 22 August 2026.

\end{thebibliography}

\end{document}
